\section{The Imaginary Bounce}

The phase face asks how orientation is carried. It does not ask the grammar
to build a new hidden object. This distinction matters here more than
anywhere else in the chapter. The quarter-turn can be read as an imaginary
direction because the loop has already carried charge, strain, phase, and
sameness far enough to make orientation accountable. But the reading is still
a reading. No completed spinor object is claimed inside this chapter.

Call the quarter-turn \(-i\) if the notation helps. The notation is a tag for
the phase cost of turning out of the real count face into the orientation
face. A full sign flip is not that single turn. The sign flip is what the
grammar reads after the phase has bounced through the needed return. In the
familiar shorthand, two such quarter-turns carry the sign change:
\[
  (-i)^2 = -1 .
\]
The equation is a disciplined mnemonic, not a new primitive. It records how
the phase reading explains why a sign can be reversed by orientation cost.

The bounce begins from the flat baseline. Over that baseline the charged loop
reads the negative matter sign. The path has returned, the count is closed,
and the native orientation faces the reader as \(-1\). A tilted comparison
does not create a second substance. It changes the facing of the loop. When
the phase reading crosses the quarter-turn and returns through the
appropriate closure, the same carried activity can be read with the opposite
orientation. The matter/antimatter split is the sign of that baseline change.

The word imaginary is therefore not a license for fantasy. It marks a
reading that is orthogonal to the count face. The real count can say how many
loops have closed. The strain face can say how the response bends. The phase
face can say how the closed loop is oriented relative to a comparison that
may have turned out of the count line. Imaginary means that the orientation
is being read in the phase direction, not that an unmeasured entity has been
added to the ledger.

The grammar also bounds the bounce. A quarter-turn alone is a phase cost. A
half-turn reads a sign change. Four quarter-turns restore the original
orientation. This is why spin one-half examples are suggestive. A \(2\pi\)
turn may not restore the sign, while a \(4\pi\) turn does. The example shows
how orientation can require a double return. But the chapter must keep its
claim narrower: the phase face reads a quarter-turn cost whose repeated
return accounts for the sign split. It does not derive a full spinor
structure.

The Aharonov-Bohm example gives another useful lantern. Local field readings
can be silent along separate branches, while the closed comparison returns
with a phase shift. The lesson for this chapter is not a detailed theory of
electromagnetism. It is the grammar of phase: transported orientation can
become readable only when the loop closes, and the closure can reveal a
global residue that no open segment could carry as charge.

The Mach-Zehnder fold makes the same point in a finite boundary form. Two
paths can remain admissible until recombination forces a common ledger. If
their carried phases agree, the fold identifies them. If they oppose, the
fold refuses a single consistent record. The branch did not by itself decide
the outcome. The boundary comparison read the phase. That is the pattern the
imaginary bounce uses: the bounce is read at closure.

This reading also protects the positron. A positive sign is not available
because notation has both \(+1\) and \(-1\). The positive sign is available
only when the phase/baseline relation permits it. Flat closure reads the
native negative sign. Tilted closure can read the positive sign. The
quarter-turn language explains how orientation cost makes the flip
intelligible, but the sign remains a loop reading under a specified baseline.

The phase face therefore has two obligations. It must be strong enough to
account for the sign bounce. It must be weak enough not to claim more than
the grammar has earned. Strong enough means that the phase reading carries
the orientation debt that distinguishes flat from tilted closure. Weak
enough means that the chapter marks the imaginary axis as an interpretive
reading on real carriers. The grammar reads the bounce; it does not pretend
to have finished every algebraic object that later physics may use.

Now the electron can be brought to the superconducting boundary. The
electron is the native negative loop reading. The Cooper pair asks what
happens when charge is admitted not as a solitary open path but as paired
closed transport through a boundary that cancels internal antisymmetric
strain. The phase bounce does not disappear at that boundary. It becomes the
cost of keeping the pair coherent.
